% !PRINTED 90
\weberpage{90}
%<<CONT

Hiermit aber ist alles ersch\"opft, was wir von $m$ fordern m\"ussen. Denn
wenn diese Bedingung f) befriedigt ist, so brauchen wir z.\,B. nur die
$e_{1,h_1}, e_{2,h_2}, \dots e_{\varrho,h_\varrho}$ durch $p$ untheilbar,
die \"ubrigen $e$ durch $p$ theilbar anzunehmen; dann ist die Determinante

\[
\begin{vmatrix}
e_{1,h_1}\,i_{1,h_1}, & \dots & e_{1,h_\varrho}\,i_{1,h_\varrho} \\
\cdot\ \cdot\ \cdot & \cdots & \cdot\ \cdot\ \cdot \\
e_{\varrho,h_1}\,i_{\varrho,h_1}, & \dots & e_{\varrho,h_\varrho}\,i_{\varrho,h_\varrho}
\end{vmatrix}
\equiv e_{1,h_1}\,e_{2,h_2} \dots e_{\varrho,h_\varrho} \pmod{p}
\]

also sicher durch $p$ untheilbar.

Und bei dieser Annahme kann auch noch die Bedingung 6. befriedigt werden.
Um dies einzusehen, bezeichne man mit $p^\pi$ irgend eine der Invarianten,
etwa $i_1$, und mit $c_h$ den Indexmodul, der nach f) durch $p^\pi$
theilbar ist, so dass

\weq{25}{c_h = p^\pi c_{1,h}}

ist. Ist dann $q_h$ der zum Indexmodul $c_h$ geh\"orige Primtheiler von
$m$, so ist $p$ entweder gleich $q_h$ oder ein Theiler von $q_h - 1$. Nach
der zuletzt gemachten Annahme ist $e_{1,h}$ nicht durch $p$ theilbar,
w\"ahrend $e_{2,h}, e_{3,h}, \dots e_{\nu,h}$ durch $p$ theilbar sind. Es
sollen dann die Producte

\[
c_{1,h}\,e_{1,h}, \quad c_{2,h}\,e_{2,h}, \; \dots \; c_{\nu,h}\,e_{\nu,h}
\]

nicht alle durch $q_h$ oder durch $q_h - 1$ theilbar sein.

Ist zun\"achst $q_h$ ein einfacher Theiler von $m$, so ist $c_h = q_h -
1$, und wegen (25) ist $c_{1,h}$ nicht durch $q_h - 1$ theilbar. Man kann
also $e_{1,h}$ noch so annehmen, dass $c_{1,h}\,e_{1,h}$ nicht durch
$q_h - 1$ theilbar ist. Ist aber $q_h$ ein $\varkappa$-facher Theiler von
$m$ und $\varkappa > 1$, und ist $q_h = p$, so ist, wenn $p$ ungerade
ist, $c_h = p^{\varkappa - 1}(p-1)$, und nach der Voraussetzung b)
$\pi \leqq \varkappa - 1$. Wenn also $c_h$ durch $p^\pi$ theilbar sein
soll, so muss $\pi = \varkappa - 1$ sein, und (25) ergiebt $c_{1,h} =
p - 1$. Ist aber $p = 2$, so ist $c_h = 2$ oder $= 2^{\varkappa - 1}$,
und es muss also nach der Voraussetzung d) $\pi = 1$ oder $= \varkappa -
2$ sein, und es ist $c_{1,h} = 1$. Man kann also
% [sic: Weber prints pi <= kappa-1; hypothesis b) gives pi >= kappa-1, which is what makes divisibility force equality.]
% [sic: Weber prints c_h=2^{kappa-1} in the second 2-adic case; the index modulus is 2^{kappa-2}, as stated in Theorem 4, and that value is required for c_{1,h}=1.]

\noindent\rule{4cm}{0.4pt}

\footnotesize
\gesperrt{Dirichlet} mit Anwendung der Integralrechnung bewiesen ist. (Abhandlungen
der Berliner Akademie 1837; \gesperrt{Dirichlet's Werke}, Bd.\,I, Nr.\,21;
\gesperrt{Dirichlet-Dedekind}, Vorlesungen \"uber Zahlentheorie, Supplement VI;
\gesperrt{Bachmann}, Analytische Zahlentheorie, Abschn.\,IV). Behalten wir die oben
benutzte Bezeichnung bei, so kann man, um f\"ur ein gegebenes $p$ der
Forderung f) zu gen\"ugen, $p^{\pi+1}$ als Factor in $m$ aufnehmen. Dann
aber m\"ussen, wenn $\varrho > 1$ ist, noch $\varrho - 1$ Primfactoren
$q$ in $m$ vorkommen, die an Congruenzen von der Form $q \equiv 1
\pmod{i_1} \dots$ gebunden sind.
\normalsize
%CONT>>
